\section{Discusión y problemas abiertos}\label{sec:discusion}
La comparación de dos cortes da la expansión relativa
\(Q(n)/Q_{\rm pred}(n)=1+\cerr/n+o(1/n)\). La ecuación~\eqref{eq:main} determina \(\cerr\) mediante
el perfil y los tres términos de frontera. Quedan dos cuestiones distintas.

La primera es certificar el signo y las cifras de \(\cerr\). El intervalo global obtenido contiene
cero. Estrecharlo exige
controlar los valores nodales completos y la media con signo \(J\). Una vía sería demostrar un
resto uniforme para la cuadratura espacial de las dieciséis contribuciones físicas. Ese resto
tendría que controlar las exclusiones diagonales y los núcleos con resolvente, en particular la
contracción \(Q_2\). En la vía del continuo, un encierro que conserve la dependencia entre
contribuciones con signo tendría que mejorar el semiancho observado en el décimo nodo GL20.

La segunda concierne al resto del conteo. Después de restar \(\cerr/n\) a \(R_Q(n)\), queda un error
\(o(1/n)\). Se desconoce su orden y si en él persiste la oscilación de periodo cuatro que se
observa a tamaños pequeños. El control puntual de la llegada
y la estabilidad con vetos arbitrarios son estimaciones más fuertes que las usadas aquí.
